\section{宇宙学}\label{sec:cosmology}

宇宙学研究宇宙的整体行为. 实验观测 (主要是宇宙背景辐射的观测) 表明, 空间在地球所在位置是各向同性的. 假定 \textbf{Copernicus 原理}成立, 即地球在宇宙中的位置并无任何特殊之处, 我们取一个各向同性 (因而具有常曲率) 的三维 Riemann 流形 $(\Sigma,h)$ 作为空间模型. 总能在 $\Sigma$ 上找到局部坐标 $(r,\theta,\varphi)$ 使得
\[
  h=a^2\left(\frac{1}{1-kr^2}\,\mathrm{d}r\otimes\mathrm{d}r
     +r^2\,\mathrm{d}\theta\otimes\mathrm{d}\theta
     +r^2\sin^2\theta\,\mathrm{d}\varphi\otimes\mathrm{d}\varphi\right),
\]
其中 $a>0$ 为空间的\textquotedblleft 半径\textquotedblright, 而 $k=-1,0,1$ 分别对应于曲率为负, 为零或为正 [参见 Exercise~\ref{ex:cosmology}(\ref{ex:constant-curvature-of-flrw-spatial-metric})]. 考虑到空间的\textquotedblleft 半径\textquotedblright 可能随时间变化, 取 $(M,g)$ 为宇宙模型, 其中 $M=\mathbb{R}\times\Sigma$, 且
\begin{align*}
  g&=-\mathrm{d}t\otimes\mathrm{d}t+a^2(t)
       \left(\frac{1}{1-kr^2}\,\mathrm{d}r\otimes\mathrm{d}r
          +r^2\,\mathrm{d}\theta\otimes\mathrm{d}\theta\right.\\
   &\qquad\qquad\qquad\qquad\qquad
       \left.{}+r^2\sin^2\theta\,\mathrm{d}\varphi\otimes\mathrm{d}\varphi\right).
\end{align*}
这些就是所谓的 \textbf{Friedmann–Lemaître–Robertson–Walker (FLRW) 宇宙学模型}.

容易计算度量 $g$ 的 Ricci 曲率. 有
\[
  g=-\omega^0\otimes\omega^0+\omega^r\otimes\omega^r
    +\omega^\theta\otimes\omega^\theta+\omega^\varphi\otimes\omega^\varphi,
\]
其中
\begin{align*}
  \omega^0&=\mathrm{d}t;\\
  \omega^r&=a(t)(1-kr^2)^{-1/2}\,\mathrm{d}r;\\
  \omega^\theta&=a(t)r\,\mathrm{d}\theta;\\
  \omega^\varphi&=a(t)r\sin\theta\,\mathrm{d}\varphi,
\end{align*}
故 $\{\omega^0,\omega^r,\omega^\theta,\omega^\varphi\}$ 是标准正交余标架. 第一结构方程给出
\begin{align*}
  \omega_r^0&=\omega_0^r=\dot a(1-kr^2)^{-1/2}\,\mathrm{d}r;\\
  \omega_\theta^0&=\omega_0^\theta=\dot a r\,\mathrm{d}\theta;\\
  \omega_\varphi^0&=\omega_0^\varphi=\dot a r\sin\theta\,\mathrm{d}\varphi;\\
  \omega_r^\theta&=-\omega_\theta^r=(1-kr^2)^{1/2}\,\mathrm{d}\theta;\\
  \omega_r^\varphi&=-\omega_\varphi^r=(1-kr^2)^{1/2}\sin\theta\,\mathrm{d}\varphi;\\
  \omega_\theta^\varphi&=-\omega_\varphi^\theta=\cos\theta\,\mathrm{d}\varphi.
\end{align*}

用第二结构方程可以求得曲率形式:
\begin{align*}
  \Omega_r^0&=\Omega_0^r=\frac{\ddot a}{a}\,\omega^0\wedge\omega^r;\\
  \Omega_\theta^0&=\Omega_0^\theta=\frac{\ddot a}{a}\,\omega^0\wedge\omega^\theta;\\
  \Omega_\varphi^0&=\Omega_0^\varphi=\frac{\ddot a}{a}\,\omega^0\wedge\omega^\varphi;\\
  \Omega_r^\theta&=-\Omega_\theta^r
     =\left(\frac{k}{a^2}+\frac{\dot a^2}{a^2}\right)\omega^\theta\wedge\omega^r;\\
  \Omega_r^\varphi&=-\Omega_\varphi^r
     =\left(\frac{k}{a^2}+\frac{\dot a^2}{a^2}\right)\omega^\varphi\wedge\omega^r;\\
  \Omega_\theta^\varphi&=-\Omega_\varphi^\theta
     =\left(\frac{k}{a^2}+\frac{\dot a^2}{a^2}\right)\omega^\varphi\wedge\omega^\theta.
\end{align*}

标准正交标架下的曲率张量分量可从曲率形式中读出, 并进而用来计算同一标架下的 Ricci 曲率张量 $\operatorname{Ric}$ 的分量. $\operatorname{Ric}$ 在此标架下的非零分量为
\begin{align*}
  R_{00}&=-\frac{3\ddot a}{a};\\
  R_{rr}&=R_{\theta\theta}=R_{\varphi\varphi}
       =\frac{\ddot a}{a}+\frac{2\dot a^2}{a^2}+\frac{2k}{a^2}.
\end{align*}
在极大的尺度上, 预期星系和星系团的行为就像无压流体的粒子, 我们取这样的流体作为物质模型. 由各向同性, 星系的平均空间运动必为零, 因而其单位速度向量场必为 $\frac{\partial}{\partial t}$ (对应于 1-形式 $-\mathrm{d}t$). 因此 Einstein 场方程为
\[
  \operatorname{Ric}=4\pi\rho(2\,\mathrm{d}t\otimes\mathrm{d}t+g),
\]
这等价于常微分方程组
\[
  \begin{cases}
    \displaystyle-\frac{3\ddot a}{a}=4\pi\rho\\[.5ex]
    \displaystyle\frac{\ddot a}{a}+\frac{2\dot a^2}{a^2}+\frac{2k}{a^2}=4\pi\rho
  \end{cases}
  \iff
  \begin{cases}
    \displaystyle\ddot a+\frac{\dot a^2}{2a}+\frac{k}{2a}=0\\[.5ex]
    \displaystyle\rho=-\frac{3\ddot a}{4\pi a}
  \end{cases}.
\]
第一个方程可确定函数 $a(t)$, 第二个方程给出 $\rho$ (特别地, $\rho$ 必定只是时间坐标 $t$ 的函数; 这意味着在宇宙学尺度上, 物质平均密度在空间中恒定). 另一方面, 量
\[
  \frac{4\pi\rho a^3}{3}=-\ddot a a^2
\]
是常量, 因为
\[
  \frac{\mathrm{d}}{\mathrm{d}t}(-\ddot a a^2)
    =\frac{\mathrm{d}}{\mathrm{d}t}\left(\frac{a\dot a^2}{2}+\frac{ka}{2}\right)
    =a\dot a\ddot a+\frac{\dot a^3}{2}+\frac{k\dot a}{2}=0.
\]
因此有
\[
  \ddot a=-\frac{\alpha}{a^2}
\]
其中 $\alpha$ 为某个积分常数 (取 $\alpha>0$, 使 $\rho>0$). 代入 $a(t)$ 的方程, 得到一阶常微分方程
\[
  \frac{\dot a^2}{2}-\frac{\alpha}{a}=-\frac{k}{2}.
\]
可由此证明 $a(t)$ 有界当且仅当 $k=1$ [参见 Exercise~\ref{ex:cosmology}(\ref{ex:flrw-scale-factor-evolution})]. 而且在所有情形中, 对某个 $t$ 值, $a(t)$ 都会变为零 (从而 $\dot a(t)$, $\ddot a(t)$ 和 $\rho(t)$ 发散), 通常取这一时刻为 $t=0$. 这一奇点称为定义在 $t>0$ 上的解的\phantomsection\label{idx:big-bang}\textbf{大爆炸}. 曾有人认为, 它是 FLRW 模型高度对称的结果. 然而, Hawking 和 Penrose 证明, 大爆炸实际上是宇宙学模型的一种一般特征 (参见第~\RGref{sec:hawking-singularity-theorem}~节和第~\RGref{sec:penrose-singularity-theorem}~节).

函数
\[
  H(t)=\frac{\dot a}{a}
\]
被称作 \textbf{Hubble 常数} (这样的名称有些令人困惑). 由上面的方程容易看出,
\[
  H^2+\frac{k}{a^2}=\frac{8\pi}{3}\rho.
\]
所以在这些模型中, $k=-1$, $k=0$ 或 $k=1$ 分别对应于宇宙的平均密度 $\rho$ 小于, 等于或大于所谓的\textbf{临界密度}
\[
  \rho_c=\frac{3H^2}{8\pi}.
\]

这些模型曾长期是宇宙学的标准模型. 然而目前认为情形要稍复杂一些 [参见 Exercise~\ref{ex:cosmology}(\ref{ex:flrw-with-cosmological-constant})].

\begin{RiemannExercise}\label{ex:cosmology}
\begin{enumerate}[label=(\arabic*),ref=\arabic*,leftmargin=*]
\item\label{ex:constant-curvature-of-flrw-spatial-metric}
证明在局部坐标 $(r,\theta,\varphi)$ 下由
\[
  h=a^2\left(\frac{1}{1-kr^2}\,\mathrm{d}r\otimes\mathrm{d}r
     +r^2\,\mathrm{d}\theta\otimes\mathrm{d}\theta
     +r^2\sin^2\theta\,\mathrm{d}\varphi\otimes\mathrm{d}\varphi\right)
\]
给出的 Riemann 度量 $h$ 具有常曲率 $K=\frac{k}{a^2}$.
\item\label{ex:flrw-comoving-observers-are-geodesic}
FLRW 模型中星系和星系群的运动是 $\frac{\partial}{\partial t}$ 的积分曲线. 证明这些曲线是类时测地线, 且时间坐标 $t$ 是此类观察者的固有时.
\item\label{ex:hubble-law-and-cosmological-redshift}
在一个 FLRW 模型中考虑两个星系, 可以假定它们的空间位置分别是 $r=0$ 和 $(r,\theta,\varphi)=(r_1,\theta_1,\varphi_1)$. 证明:
\begin{enumerate}[label=(\alph*),leftmargin=*]
\item 两个星系沿 $t$ 为常数的空间 Riemann 流形的空间距离 $d(t)$ 满足 \textbf{Hubble 定律}
\[
  \dot d=Hd,
\]
其中 $H=\dot a/a$ 为 Hubble 常数;
\item 将第一个星系与第二个星系连接起来的 (重新参数化的) 零模测地线族可写成
\[
  (t,r,\theta,\varphi)=(t(r,t_0),r,\theta_1,\varphi_1)\qquad(0<r<r_1),
\]
其中 $t(r,t_0)$ 为方程组
\[
  \begin{cases}
    \dfrac{\mathrm{d}t}{\mathrm{d}r}=\dfrac{a(t)}{\sqrt{1-kr^2}}\\[.5ex]
    t(0,t_0)=t_0
  \end{cases}
\]
的解;
\item $\dfrac{\partial t}{\partial t_0}(r_1,t_0)=\dfrac{a(t_1)}{a(t_0)}$, 其中 $t_1=t(r_1,t_0)$.

(\textbf{注}: 这意味着, 第一个星系以周期 $T$ 发出的光, 被第二个星系测得的周期为 $T'=\frac{a(t_1)}{a(t_0)}T$.)
\end{enumerate}
\item\label{ex:flrw-scale-factor-evolution}
回顾, FLRW 模型中空间的\textquotedblleft 半径\textquotedblright $a(t)$ 按照常微分方程
\[
  \frac{\dot a^2}{2}-\frac{\alpha}{a}=-\frac{k}{2}
  \implies\ddot a=-\frac{\alpha}{a^2}
\]
演化. 证明:
\begin{enumerate}[label=(\alph*),leftmargin=*]
\item $a(t)$ 在有限时间变为零 (假定它发生在 $t=0$);
\item 若 $k=-1$ 或 $k=0$, 则解可以延拓到所有 $t>0$;
\item 若 $k=1$, 则解不能延拓到 $t=2\pi\alpha$ 之后 (\phantomsection\label{idx:big-crunch}\textbf{大坍缩});
\item 若 $k=1$, 则无论观察者运动得多快, 都不可能绕宇宙一周;
\item 解可以用参数方程给出:
\begin{enumerate}[label=(\roman*),leftmargin=*]
\item $k=1$:
\[
  \begin{cases}
    a=\alpha(1-\cos u)\\
    t=\alpha(u-\sin u)
  \end{cases};
\]
\item $k=0$:
\[
  \begin{cases}
    a=\frac{\alpha}{2}u^2\\
    t=\frac{\alpha}{6}u^3
  \end{cases};
\]
\item $k=-1$:
\[
  \begin{cases}
    a=\alpha(\cosh u-1)\\
    t=\alpha(\sinh u-u)
  \end{cases}.
\]
\end{enumerate}
\end{enumerate}
\item\label{ex:closed-flrw-embedding-in-minkowski-spacetime}
证明 $k=1$ 的 FLRW 模型, 与五维 Minkowski 时空 $(\mathbb{R}^5,g)$ 中方程为
\[
  \sqrt{x^2+y^2+z^2+w^2}=2\alpha-\frac{t^2}{8\alpha}
\]
的超曲面等距, 其中度量为
\[
  g=-\mathrm{d}t\otimes\mathrm{d}t+\mathrm{d}x\otimes\mathrm{d}x
    +\mathrm{d}y\otimes\mathrm{d}y+\mathrm{d}z\otimes\mathrm{d}z+\mathrm{d}w\otimes\mathrm{d}w.
\]
\item\label{ex:flrw-model-for-gravitational-collapse}
(\emph{一个坍缩模型}) 证明若
\[
  \begin{cases}
    m=\alpha r_0^3\\
    E^2-1=-kr_0^2
  \end{cases},
\]
则 FLRW 模型中球壳 $r=r_0$ 的半径随固有时变化的方式, 与质量参数为 $m$ 的 Schwarzschild 时空中以能量参数 $E$ 作径向自由下落运动的球壳的半径变化方式完全相同 [参见 Exercise~\ref{ex:schwarzschild-solution}(\ref{ex:schwarzschild-orbits-and-light-deflection})]. 因此, 可沿球壳运动确定的三维超曲面将这两个时空匹配, 得到一个物质坍缩的模型. 你能给出这个模型的物理解释吗?
\item\label{ex:flrw-with-cosmological-constant}
证明若允许存在\textbf{宇宙学常数} $\Lambda\in\mathbb{R}$, 即 Einstein 方程具有形式
\[
  \operatorname{Ric}=4\pi\rho(2\nu\otimes\nu+g)+\Lambda g,
\]
则 FLRW 模型的方程变为
\[
  \begin{cases}
    \dfrac{\dot a^2}{2}-\dfrac{\alpha}{a}-\dfrac{\Lambda}{6}a^2=-\dfrac{k}{2}\\[.6ex]
    \dfrac{4\pi}{3}a^3\rho=\alpha
  \end{cases}.
\]
分析函数 $a(t)$ 的可能行为.

(\textbf{注}: 目前认为, 确实存在正的宇宙学常数, 也称为\textbf{暗能量}. 实验观测所支持的模型似乎是 $k=0$, $\Lambda>0$.)
\item\label{ex:einstein-and-de-sitter-universe-embeddings}
考虑度量为
\[
  g=-\mathrm{d}t\otimes\mathrm{d}t+\mathrm{d}x\otimes\mathrm{d}x
    +\mathrm{d}y\otimes\mathrm{d}y+\mathrm{d}z\otimes\mathrm{d}z+\mathrm{d}w\otimes\mathrm{d}w
\]
的五维 Minkowski 时空 $(\mathbb{R}^5,g)$. 证明以下各超曲面上的诱导度量确定了具有所列参数的 FLRW 模型.
\begin{enumerate}[label=(\alph*),leftmargin=*]
\item \textbf{Einstein 宇宙}: 方程为
\[
  x^2+y^2+z^2+w^2=\frac{1}{\Lambda}
\]
的\textquotedblleft 柱面\textquotedblright 满足 $k=1$, $\Lambda>0$ 和 $\rho=\frac{\Lambda}{4\pi}$.
\item \textbf{de Sitter 宇宙}: 方程为
\[
  -t^2+x^2+y^2+z^2+w^2=\frac{3}{\Lambda}
\]
的\textquotedblleft 球面\textquotedblright 满足 $k=1$, $\Lambda>0$ 和 $\rho=0$.
\end{enumerate}
\item\label{ex:flrw-brightness-and-redshift}
若光信号以周期 $T$ 发出, 以周期 $T'$ 接收, 则称它经历了\textbf{红移}
\[
  z=\frac{T'}{T}-1
\]
[所以在 Doppler 效应的情形下, 对小速度有 $z\simeq v$, 参见 Exercise~\ref{ex:special-relativity}(\ref{ex:relativistic-doppler-effect})]. 若光由位于 $r=0$ 的星系在时刻 $t=t_0$ 发出, 由位于 $r=r_1$ 的星系在时刻 $t=t_1$ 接收, 则它的红移为
\[
  z=\frac{a(t_1)}{a(t_0)}-1
\]
[参见 Exercise~\ref{ex:cosmology}(\ref{ex:hubble-law-and-cosmological-redshift})]. 这束光散布在半径为 $R=a(t_1)r_1$ 的球面上, 因而其亮度与 $R^2$ 成反比. 对下列 FLRW 模型, 求 $R$ 作为 $z$ 的函数:
\begin{enumerate}[label=(\alph*),leftmargin=*]
\item \textbf{Milne 宇宙} ($k=-1$, $\alpha=\Lambda=0$), 其中 $a(t)=t$;
\item \textbf{平直 de Sitter 宇宙} ($k=\alpha=0$, $\Lambda=3H^2$), 其中 $a(t)=e^{Ht}$;
\item \textbf{Einstein–de Sitter 宇宙} ($k=\Lambda=0$, $\alpha=2/(9t_1^2)$), 其中 $a(t)=(t/t_1)^{2/3}$.
\end{enumerate}
(\textbf{注}: 遥远星系的亮度还会按因子 $(1+z)^2$ 进一步减小, 因为接收时每个光子的频率, 从而能量, 都变为发出时的 $1/(1+z)$, 而光子的探测速率也变为发射速率的 $1/(1+z)$. 作此修正后, 对于光度已知的星系, 可以从观测亮度推得 $R$, 并把曲线 $R=R(z)$ 最符合观测的 FLRW 模型选为正确模型.)
\end{enumerate}
\end{RiemannExercise}
